% !TEX root = ../main.tex

\section{Extension 1: Dynamic Misspecification (3-4 pages)}

Ultimately, the baseline design establishes a clear reference under correct specification. The equal-lag VAR is the more efficient point estimator, while the local projection offers substantially more reliable inference. The open question remains whether the robustness argument for local projections prevails once the theoretical assumption of correct specification is relaxed. The subsequent extensions address this directly. We first introduce controlled dynamic misspecification to target the bias channel where the complexity adjustment theoretically binds, followed by functional misspecification to isolate non-linear effects in the conditional mean.

\subsection{Simulation Design}

Extension 1 examines whether the finite-sample MSE ranking between $\text{LP}(4)$ and the complexity-adjusted VAR sweep is robust to dynamic misspecification of the Wold representation. The estimators, the structural estimand, the recursive Cholesky identification, and all performance measures are inherited unchanged from the baseline. Only the data-generating process is modified. Specifically, the baseline $\text{VAR}(4)$ is augmented by a first-order moving-average component, yielding a $\text{VARMA}(4,1)$ DGP:

\begin{equation}
    y_t = A_1 y_{t-1} + A_2 y_{t-2} + A_3 y_{t-3} + A_4 y_{t-4} 
    + u_t + \Theta u_{t-1}, \quad u_t = B\varepsilon_t, \quad 
    \varepsilon_t \overset{\text{iid}}{\sim} \mathcal{N}(\mathbf{0}, I_2)
\end{equation}

where $\Theta = \theta I_2$ is a scalar-scaled MA impact matrix acting on the lagged reduced-form innovation $u_{t-1}$, with misspecification strength $\theta \in \{0.1, 0.3, 0.6\}$. The case $\theta = 0$ nests back to the baseline $\text{VAR}(4)$, ensuring continuity across designs. The autoregressive matrices $A_1, \ldots, A_4$ and the impact matrix $B$ are inherited from the baseline DGP, so that all variation across the Extension 1 scenarios is attributable exclusively to the MA component.

The $\text{VARMA}(4,1)$ DGP is invertible for all $\theta \in \{0.1, 0.3, 0.6\}$, meaning the roots of $\det(I_2 + \Theta z)$ lie outside the unit circle. The process therefore admits an infinite-order VAR representation. This preserves the theoretical motivation of the complexity-adjusted VAR sweep. Higher-order $\text{VAR}(q)$ specifications progressively approximate the true dynamics as $q \to \infty$, though at the cost of increasing parameter proliferation in finite samples. The central question is therefore whether $\text{LP}(4)$, which remains consistent under this approximable dynamic misspecification by virtue of its zero asymptotic bias \parencite{olea2025local}, retains its MSE advantage over the equal-lag $\text{VAR}(4)$, and where on the complexity-adjusted VAR frontier it sits relative to the baseline.

The design dimensions partially follow the baseline. Persistence is varied across $\rho \in \{0.5, 0.7, 0.95\}$ via the spectral radius of the companion matrix. The choice of $\rho$ is critical here. The degree to which the MA component distorts the finite-sample VAR approximation depends directly on the persistence of the underlying process. At high $\rho$, shocks decay slowly and the MA term remains influential at longer horizons, amplifying the VAR approximation bias. To isolate the marginal contribution of MA misspecification from the finite-sample channels already documented in the baseline, the sample size is fixed at $T = 250$ observations, generated after a 200-period burn-in. The misspecification strength $\theta$ is thus varied jointly with persistence, yielding a two-way design.

The structural estimand retains the baseline form $\theta_h = \Psi_h B e_1$, but is now defined with respect to the true VARMA Wold matrices $\Psi_h^{\mathrm{VARMA}}$. These solve $A(L)\Psi(L) = I_2 + \Theta L$ and are computed analytically prior to each simulation run via
\begin{equation}
    \Psi_0 = I_2, \qquad 
    \Psi_h = \sum_{i=1}^{\min(h,4)} A_i \Psi_{h-i} + \Theta_h, 
    \qquad \Theta_h = \begin{cases} \Theta & h = 1 \\ \mathbf{0} & h > 1 \end{cases}
\end{equation}
so that the additive MA term $\Theta_h$ contributes only at horizon $h = 1$. Both estimators, $\text{LP}(4)$ and the $\text{VAR}(q)$ sweep, are applied without modification to the VARMA-generated data, meaning the misspecification is fully absorbed into the finite-sample performance of each estimator rather than accommodated by the estimation procedure. All performance measures and the $B = 5{,}000$ replication count are identical to the baseline.

% ----------------------------------- True IRF
\begin{figure}[htbp]
\centering
\caption{True Impulse Response Functions Across MA Misspecification}
\label{fig:DYN_TRUE_IRF}


\begin{subfigure}[t]{0.48\textwidth}
    \centering
    \caption{\(\rho = 0.5\)}
    \includegraphics[width=\textwidth]{DYN_MISSPEC_IRF_rho=0.5_crop.png}
    \label{fig:DYN_TRUE_IRF_LOW}
\end{subfigure}
\hfill
\begin{subfigure}[t]{0.48\textwidth}
    \centering
    \caption{\(\rho = 0.7\)}
    \includegraphics[width=\textwidth]{DYN_MISSPEC_IRF_rho=0.7_crop.png}
    \label{fig:DYN_TRUE_IRF_MID}
\end{subfigure}

\begin{subfigure}[t]{0.48\textwidth}
    \centering
    \caption{\(\rho = 0.95\)}
    \includegraphics[width=\textwidth]{DYN_MISSPEC_IRF_rho=0.95_crop.png}
    \label{fig:DYN_TRUE_IRF_HIGH}
\end{subfigure}

\begin{minipage}{0.9\textwidth}
\small
\textit{Note.} Each panel reports the true structural impulse response function of variables one and two to a one-standard-deviation shock in the first structural innovation, traced over horizons 
\(h = 1, \ldots, 20\). Panels correspond to persistence regimes \(\rho \in \{0.5, 0.7, 0.95\}\). Blue line represents $y_1 \leftarrow \epsilon_1 ($\textit{estimand} $ \theta_h)$, while grey line shows $y_2 \leftarrow \epsilon_1$.
\end{minipage}
\end{figure}

\FloatBarrier

\subsection{Simulation Results}


Figure \ref{fig:DYN_MISSPEC_COMPLEXITY} places the fixed $\text{LP}(4)$ estimator on the VAR-order frontier. Each panel plots the horizon-averaged root mean squared error of the structural impulse response as a function of the VAR order $q$, with one solid curve for each misspecification strength $\theta$ and the matching $\text{LP}(4)$ value drawn as a horizontal dashed line in the corresponding color. As $\theta$ increases, the entire frontier shifts upward, and the $\text{LP}(4)$ reference line rises accordingly. The minimum of the frontier consistently occurs at a low lag order, typically between $\text{VAR}(1)$ and $\text{VAR}(3)$.

Across all persistence regimes, the equal-lag $\text{VAR}(4)$ and the frontier minimum lie strictly below the $\text{LP}(4)$ reference line for every value of $\theta$. The $\text{LP}(4)$ line intersects the rising VAR curves only at higher orders, with the intersection points generally located between $q = 16$ and $q = 18$. Furthermore, the performance of the lowest-order models deteriorates noticeably as $\theta$ and $\rho$ increase. At $\rho = 0.95$ with $\theta = 0.6$, the curves for $\text{VAR}(1)$ and $\text{VAR}(2)$ sit substantially higher than the heavily parameterized $\text{VAR}(23)$ specification, producing a jagged shape at the low-order end of the frontier where $\text{VAR}(2)$ yields the highest average error.
% ----------------------------------- Complexity Frontier
\begin{figure}[htbp]
\centering
\caption{Complexity Frontier under Dynamic Misspecification}
\label{fig:DYN_MISSPEC_COMPLEXITY}

\begin{subfigure}[t]{0.48\textwidth}
    \centering
    \caption{\(\rho = 0.5\)}
    \includegraphics[width=\textwidth]{DYN_MISSPEC_COMPLEXITY_FRONTIER_rho=0.5_T=250.png}
    \label{fig:DYN_MISSPEC_COMPLEXITY_rho05}
\end{subfigure}
\hfill
\begin{subfigure}[t]{0.48\textwidth}
    \centering
    \caption{\(\rho = 0.7\)}
    \includegraphics[width=\textwidth]{DYN_MISSPEC_COMPLEXITY_FRONTIER_rho=0.7_T=250.png}
    \label{fig:DYN_MISSPEC_COMPLEXITY_rho07}
\end{subfigure}

\begin{subfigure}[t]{0.48\textwidth}
    \centering
    \caption{\(\rho = 0.95\)}
    \includegraphics[width=\textwidth]{DYN_MISSPEC_COMPLEXITY_FRONTIER_rho=0.95_T=250.png}
    \label{fig:DYN_MISSPEC_COMPLEXITY_rho095}
\end{subfigure}

\begin{minipage}{0.9\textwidth}
\small
\textit{Note.} Horizon-averaged RMSE of the structural IRF estimate against VAR order \(q = 1,\ldots,23\) (solid), with LP(4) as a same-colour horizontal dashed reference; the vertical dotted line marks the equal-lag VAR(4). One colour per MA strength: purple \(\theta=0\) (baseline), blue \(\theta=0.1\), green \(\theta=0.3\), yellow \(\theta=0.6\). \(T=250\), \(B=5{,}000\); panels \(\rho \in \{0.5,0.7,0.95\}\). As \(\theta\) grows the whole frontier and its LP(4) reference lift.
\end{minipage}

\end{figure}

Figure \ref{fig:DYN_MISSPEC_POINT_ESTIMATION} decomposes the finite-sample error at the strongest misspecification level $\theta = 0.6$ and $T = 250$ by reporting the bias, variance, mean squared error, and root mean squared error of the estimated impulse response $\hat\theta_h$ across horizons. The $\text{LP}(4)$ estimator is compared against a representative set of VAR orders.

The bias panel separates the estimators into distinct trajectories. The $\text{LP}(4)$ estimator exhibits a roughly constant negative bias that deepens as persistence rises. The well-lagged VAR specifications from $\text{VAR}(4)$ upward closely track this path. They oscillate over the first few horizons before approaching zero, placing the equal-lag $\text{VAR}(4)$ and $\text{LP}(4)$ in close proximity. The $\text{VAR}(2)$ specification deviates from this pattern, displaying a substantially larger bias that grows further as the persistence parameter increases.

In the variance panel, the estimators are ordered systematically by lag length, with variance increasing alongside the VAR order. The $\text{LP}(4)$ estimator aligns with the higher-variance VAR models. The parsimonious $\text{VAR}(2)$ specification exhibits one of the lowest variance profiles across all horizons.

The mean squared error and root mean squared error panels combine these two components. The $\text{VAR}(2)$ specification shifts to the top of the panels, exceeding all other estimators by a wide margin. The longer-lagged VAR specifications, particularly $\text{VAR}(4)$ and $\text{VAR}(8)$, maintain a lower error level than $\text{LP}(4)$ across the entire horizon window. The $\text{VAR}(4)$ model therefore maintains a root mean squared error advantage over $\text{LP}(4)$ across all persistence levels. The overall curvature of these error profiles becomes more pronounced as persistence increases.

% ----------------------------------- Point Estimates
\begin{figure}[htbp]
\centering
\caption{Point-Estimation Performance Across Persistence Regimes}
\label{fig:DYN_MISSPEC_POINT_ESTIMATION}

\begin{subfigure}[t]{0.6\textwidth}
    \centering
    \caption{\(\rho = 0.5\)}
    \includegraphics[width=\textwidth]{DYN_MISSPEC_POINT_ESTIMATION_rho=0.5_T=250_B=5000_theta=0.6.png}
    \label{fig:DYN_MISSPEC_POINT_ESTIMATION_rho05}
\end{subfigure}

\begin{subfigure}[t]{0.6\textwidth}
    \centering
    \caption{\(\rho = 0.7\)}
    \includegraphics[width=\textwidth]{DYN_MISSPEC_POINT_ESTIMATION_rho=0.7_T=250_B=5000_theta=0.6.png}
    \label{fig:DYN_MISSPEC_POINT_ESTIMATION_rho07}
\end{subfigure}

\begin{subfigure}[t]{0.6\textwidth}
    \centering
    \caption{\(\rho = 0.95\)}
    \includegraphics[width=\textwidth]{DYN_MISSPEC_POINT_ESTIMATION_rho=0.95_T=250_B=5000_theta=0.6.png}
    \label{fig:DYN_MISSPEC_POINT_ESTIMATION_rho095}
\end{subfigure}

\begin{minipage}{0.9\textwidth}
\small
\textit{Note.} Bias, variance, MSE, and RMSE of \(\hat\theta_h\) relative to the true estimand \(\theta_h\), horizon-by-horizon, at the strongest misspecification \(\theta=0.6\), \(T=250\), \(B=5{,}000\). LP(4) (red) is compared against a representative VAR\((q)\) set: VAR(2) blue, VAR(4) black, VAR(8) light green, VAR(16) dark green, VAR(23) purple. Panels \(\rho \in \{0.5,0.7,0.95\}\); the shaded band on bias is the \(95\%\) Monte Carlo standard-error band.
\end{minipage}

\end{figure}

\newpage

Figure \ref{fig:DYN_MISSPEC_COVERAGE} visualizes the empirical coverage of nominal 95 percent confidence intervals over horizons $h = 1$ to 20 at $T = 250$. The left and middle panels display the coverage for the $\text{VAR}(4)$ and $\text{LP}(4)$ models, featuring one curve per misspecification strength $\theta$. The right panel plots the $\text{VAR}(4)$ delta-method standard error alongside the empirical sampling standard deviation. Within each panel, the four $\theta$ curves are tightly clustered and lie almost perfectly on top of one another. The level and shape of these curves shift noticeably across the three persistence regimes.

For the $\text{VAR}(4)$ model, the intervals at $\rho = 0.5$ reach full coverage at longer horizons. At $\rho = 0.7$, the coverage drops slightly below the nominal target, and at $\rho = 0.95$, it deteriorates further, declining steadily as the horizon increases. The $\text{LP}(4)$ coverage remains relatively steady and near the 0.95 mark at low and intermediate persistence levels. At $\rho = 0.95$, the $\text{LP}(4)$ coverage slips to approximately 0.89, maintaining a higher level than the $\text{VAR}(4)$ specification but exhibiting a similar downward slope over longer horizons.

Varying the moving-average parameter $\theta$ produces only minor shifts within the panels. At $\rho = 0.5$, the $\text{VAR}(4)$ coverage drifts downward slightly as $\theta$ increases. At $\rho = 0.7$, the coverage improves marginally from roughly 0.89 at $\theta = 0$ to about 0.93 at $\theta = 0.6$. At $\rho = 0.95$, the coverage curves are virtually indistinguishable across all values of $\theta$.

In the right panel, the delta-method standard error and the empirical standard deviation show almost no deviation at $\rho = 0.5$ and $\rho = 0.7$. At $\rho = 0.95$, a vertical gap emerges, with the empirical dispersion sitting visibly above the analytical standard error over the short and intermediate horizons.

% ----------------------------------- Coverage
\begin{figure}[htbp]
\centering
\caption{Confidence-Interval Coverage under Dynamic Misspecification}
\label{fig:DYN_MISSPEC_COVERAGE}

\begin{subfigure}[t]{0.9\textwidth}
    \centering
    \caption{\(\rho = 0.5\)}
    \includegraphics[width=\textwidth]{DYN_MISSPEC_COVERAGE_rho=0.5_T250_B=5000.png}
    \label{fig:DYN_MISSPEC_COVERAGE_rho05}
\end{subfigure}

\begin{subfigure}[t]{0.9\textwidth}
    \centering
    \caption{\(\rho = 0.7\)}
    \includegraphics[width=\textwidth]{DYN_MISSPEC_COVERAGE_rho=0.7_T250_B=5000.png}
    \label{fig:DYN_MISSPEC_COVERAGE_rho07}
\end{subfigure}

\begin{subfigure}[t]{0.9\textwidth}
    \centering
    \caption{\(\rho = 0.95\)}
    \includegraphics[width=\textwidth]{DYN_MISSPEC_COVERAGE_rho=0.95_T250_B=5000.png}
    \label{fig:DYN_MISSPEC_COVERAGE_rho095}
\end{subfigure}

\begin{minipage}{0.9\textwidth}
\small
\textit{Note.} Empirical coverage of nominal \(95\%\) CIs by horizon (\(h=1,\ldots,20\)) at 
\(T=250\), \(B=5{,}000\). \emph{Left:} VAR(4) coverage; \emph{middle:} LP(4) coverage with 
one line per MA strength (purple \(\theta=0\) baseline, blue \(\theta=0.1\), green \(\theta=0.3\), 
yellow \(\theta=0.6\); dotted line marks the \(95\%\) reference). VAR(4) uses delta-method 
intervals, LP(4) heteroskedasticity-robust SEs. \emph{Right:} the VAR(4) delta-method SE 
(orange) against the empirical sampling SD (blue) at \(\theta=0.6\). Panels \(\rho \in 
\{0.5,0.7,0.95\}\).
\end{minipage}

\end{figure}



\FloatBarrier
\subsection{Discussion of Results \& Conclusion}

\FloatBarrier